\section{Orbits as Equivalence Classes and Conjugacy Classes}

\subsection{Equivalence Relation Induced by an Action}

\begin{theorem}
Let $*$ be a group action of a group $G$ on a set $A$, $* : G \times A \to A$. Define a relation $\sim$ on $A$ by:
\[
a \sim b \iff a = g * b \quad \text{for some } g \in G.
\]
Then $\sim$ is an equivalence relation on $A$.
\end{theorem}

\begin{proof}
We verify the three equivalence relation properties:
\begin{enumerate}[label=(\roman*)]
    \item \textbf{Reflexive:} For all $a \in A$, since $e \in G$ is the identity:
    \[
    a = e * a \implies a \sim a.
    \]
    Therefore, $\sim$ is reflexive.
    
    \item \textbf{Symmetric:} Let $a, b \in A$ such that $a \sim b$. Then there exists $g \in G$ such that:
    \[
    a = g * b.
    \]
    Applying $g^{-1} \in G$ to both sides:
    \begin{align*}
    g^{-1} * a &= g^{-1} * (g * b) \\
    &= (g^{-1} g) * b \\
    &= e * b = b.
    \end{align*}
    Since $g^{-1} \in G$, we have $b \sim a$. Therefore, $\sim$ is symmetric.
    
    \item \textbf{Transitive:} Let $a, b, c \in A$ such that $a \sim b$ and $b \sim c$. Then there exist $g_1, g_2 \in G$ such that:
    \[
    a = g_1 * b \quad \text{and} \quad b = g_2 * c.
    \]
    Substituting $b$ into the expression for $a$:
    \[
    a = g_1 * (g_2 * c) = (g_1 g_2) * c.
    \]
    Since $g_1 g_2 \in G$, we have $a \sim c$. Therefore, $\sim$ is transitive.
\end{enumerate}
Hence, $\sim$ is an equivalence relation on $A$.
\end{proof}

\subsection{Partition of the Set into Orbits}

Since $\sim$ is an equivalence relation on $A$, it partitions the set $A$ into its equivalence classes. For any $a \in A$, the equivalence class of $a$ is:
\begin{align*}
\operatorname{cl}(a) &= \{b \in A \mid a \sim b\} \\
&= \{b \in A \mid b = g * a \text{ for some } g \in G\} \\
&= \{g * a \mid g \in G\} \\
&= G^a \quad \text{(the orbit of $a$ under $G$)}.
\end{align*}
Consequently, the set $A$ decomposes as a disjoint union of its orbits:
\[
A = \bigcup_{a \in A} \operatorname{cl}(a) = \bigcup_{a \in A} G^a.
\]

\subsection{Conjugacy Classes}

If $G$ acts on itself ($A = G$) by conjugation, where $g * a = gag^{-1}$:
\begin{itemize}
    \item The stabilizer of $a$ is the \textbf{normalizer} (or centralizer) of $a$:
    \[
    G_a = N(a) = \{g \in G \mid gag^{-1} = a\} = \{g \in G \mid ga = ag\}.
    \]
    \item The orbit of $a$ is the \textbf{conjugacy class} of $a$:
    \[
    G^a = \operatorname{cl}(a) = \{g * a \mid g \in G\} = \{gag^{-1} \mid g \in G\}.
    \]
\end{itemize}

\begin{exercise}
Let $G$ be a group and $a, b \in G$. Then $a$ and $b$ are conjugate if there exists $c \in G$ such that $b = cac^{-1}$. Prove that this conjugacy relation $\sim$ is an equivalence relation on $G$.
\end{exercise}

\begin{proof}[Solution]
We check the three properties:
\begin{enumerate}[label=(\roman*)]
    \item \textbf{Reflexive:} Let $a \in G$. Since $e \in G$, we have:
    \[
    a = e a e^{-1}.
    \]
    Therefore, $a \sim a$, so $\sim$ is reflexive.
    
    \item \textbf{Symmetric:} Let $a, b \in G$ such that $a \sim b$. Then there exists $c \in G$ such that:
    \[
    a = c b c^{-1}.
    \]
    Multiplying on the left by $c^{-1}$ and on the right by $c$:
    \begin{align*}
    c^{-1} a &= c^{-1}(c b c^{-1}) = b c^{-1}, \\
    c^{-1} a c &= b c^{-1} c = b e = b.
    \end{align*}
    Writing $c = (c^{-1})^{-1}$, this gives:
    \[
    c^{-1} a (c^{-1})^{-1} = b.
    \]
    Since $c^{-1} \in G$, this shows $b \sim a$. Therefore, $\sim$ is symmetric.
    
    \item \textbf{Transitive:} Let $a, b, c \in G$ such that $a \sim b$ and $b \sim c$. Then there exist $d_1, d_2 \in G$ such that:
    \[
    a = d_1 b d_1^{-1} \quad \text{and} \quad b = d_2 c d_2^{-1}.
    \]
    Substituting $b$ into the first equation:
    \begin{align*}
    a &= d_1 (d_2 c d_2^{-1}) d_1^{-1} \\
    &= d_1 d_2 c d_2^{-1} d_1^{-1} \quad \{d_2^{-1} d_1^{-1} = (d_1 d_2)^{-1}\} \\
    &= (d_1 d_2) c (d_1 d_2)^{-1}.
    \end{align*}
    Since $d_1, d_2 \in G \implies d_1 d_2 \in G$, it follows that $a \sim c$.
\end{enumerate}
Therefore, $\sim$ is transitive, and hence an equivalence relation on $G$.
\end{proof}
